\documentclass[11pt]{article}
\usepackage{booktabs}
\usepackage{graphicx}
\usepackage{amsmath}
\usepackage{hyperref}
\usepackage[margin=1in]{geometry}

\title{Methods: Fact-Based Drug-Drug Interaction Knowledge Graph Construction}
\author{}
\date{}

\begin{document}

\maketitle

\section{Knowledge Graph Construction}

\subsection{Overview}

We constructed a comprehensive drug-drug interaction (DDI) knowledge graph integrating data from multiple biomedical databases. The graph contains 45,655 nodes and 1,211,674 edges, representing 4,313 drugs and their relationships with proteins, side effects, diseases, pathways, and drug categories.

\subsection{Data Sources}

\subsubsection{DrugBank}
The primary source for drug information was DrugBank~\cite{wishart2018drugbank}, providing:
\begin{itemize}
    \item Drug identifiers (DrugBank ID, CAS number, UNII, ATC codes)
    \item Chemical properties (SMILES, InChI, molecular weight)
    \item Drug-protein interactions (targets, enzymes, carriers, transporters)
    \item Biological pathways and drug categories
    \item Pharmacogenomic data (SNP effects and adverse reactions)
    \item Drug-drug interaction descriptions
\end{itemize}

\subsubsection{SIDER}
Side effect information was obtained from SIDER 4.1~\cite{kuhn2016sider}, containing drug-side effect associations with MedDRA terminology. Drugs were matched using exact drug name matching (case-insensitive), resulting in 1,110 matched drugs and 265,238 drug-side effect associations.

\subsubsection{CTD}
Drug-disease associations were obtained from the Comparative Toxicogenomics Database (CTD)~\cite{davis2021ctd}. Matching was performed using CAS numbers (primary) and drug names (secondary), yielding 63,278 drug-disease associations across 3,041 unique diseases.

\subsection{Graph Schema}

The knowledge graph follows a heterogeneous schema with six node types and six edge types:

\textbf{Node Types:}
\begin{itemize}
    \item \texttt{Drug} (n=4,313): Primary entities with chemical and pharmacological properties
    \item \texttt{Protein} (n=3,176): Drug targets including enzymes, transporters, and carriers
    \item \texttt{SideEffect} (n=5,548): Adverse drug reactions from SIDER
    \item \texttt{Disease} (n=3,041): Drug-disease associations from CTD
    \item \texttt{Pathway} (n=25,958): Biological pathways from DrugBank
    \item \texttt{Category} (n=3,619): Drug classifications from DrugBank
\end{itemize}

\textbf{Edge Types:}
\begin{itemize}
    \item \texttt{INTERACTS\_WITH} (n=759,774): Drug-drug interactions with severity
    \item \texttt{TARGETS} (n=21,559): Drug-protein relationships
    \item \texttt{CAUSES} (n=265,238): Drug-side effect associations
    \item \texttt{ASSOCIATED\_WITH} (n=63,278): Drug-disease relationships
    \item \texttt{PARTICIPATES\_IN} (n=31,207): Drug-pathway relationships
    \item \texttt{BELONGS\_TO} (n=70,618): Drug-category relationships
\end{itemize}

\subsection{Fact-Based Construction Methodology}

Unlike similarity-based approaches, our knowledge graph uses exclusively fact-based linkage with exact identifier matching:

\begin{enumerate}
    \item \textbf{DrugBank ID Matching}: Primary matching using exact DrugBank identifiers (4,313 drugs matched)
    \item \textbf{CAS Number Matching}: Secondary matching for CTD data using Chemical Abstracts Service numbers (52,753 associations)
    \item \textbf{Drug Name Matching}: Tertiary matching using exact case-insensitive drug names (11,635 associations)
\end{enumerate}

\subsection{Provenance Tracking}

Every edge in the knowledge graph includes provenance metadata:
\begin{itemize}
    \item \texttt{source}: Data source (DrugBank, SIDER, CTD)
    \item \texttt{match\_type}: Matching method used (drugbank\_id, cas\_number, drug\_name)
    \item \texttt{match\_value}: The specific identifier used for matching
    \item \texttt{confidence}: Match confidence (exact, verified)
\end{itemize}

\section{DDI Severity Recalibration}

\subsection{Problem Statement}

The original zero-shot severity classification using BART-large-MNLI showed significant class imbalance (56.9\% Contraindicated, 0.0\% Moderate), inconsistent with clinical literature distributions.

\subsection{Empirical Keyword-Based Classification}

We developed a rule-based classifier using empirically-derived keyword weights from the DDInter database~\cite{xiong2022ddinter}:

\begin{equation}
\text{score}(d) = \sum_{k \in \text{keywords}} w_k \cdot \mathbb{1}[k \in d]
\end{equation}

where $w_k$ represents log-likelihood ratio weights derived from DDInter training data (n=26,014):

\begin{equation}
w_k = \ln\left(\frac{P(k|\text{Major})}{P(k|\text{Moderate})}\right)
\end{equation}

Key keyword weights include: \textit{prolongation} (+1.45), \textit{bleeding} (+1.03), \textit{therapeutic} (-1.20), \textit{efficacy} (-1.27).

\subsection{Percentile-Based Classification}

Severity levels are assigned using percentile thresholds optimized via grid search:
\begin{itemize}
    \item Contraindicated: score $\geq$ P96 (top 4\%)
    \item Major: score $\geq$ P92 (top 8\%)
    \item Minor: score $\leq$ P2 (bottom 2\%)
    \item Moderate: All remaining interactions
\end{itemize}

\subsection{Validation}

The classifier achieved 66.4\% exact accuracy and Cohen's $\kappa$ = +0.096 (p < 0.0001) against DDInter ground truth (n=11,150 test pairs).

\section{Output Formats}

The knowledge graph is available in two formats:
\begin{enumerate}
    \item \textbf{NetworkX Pickle}: Complete graph with all node/edge attributes and metadata
    \item \textbf{Neo4j CSV Export}: Separate CSV files for each node and edge type, compatible with Neo4j import
\end{enumerate}

\bibliographystyle{plain}
\begin{thebibliography}{9}

\bibitem{wishart2018drugbank}
Wishart DS, et al.
DrugBank 5.0: a major update to the DrugBank database for 2018.
\textit{Nucleic Acids Res}. 2018;46(D1):D1074-D1082.

\bibitem{kuhn2016sider}
Kuhn M, et al.
The SIDER database of drugs and side effects.
\textit{Nucleic Acids Res}. 2016;44(D1):D1075-D1079.

\bibitem{davis2021ctd}
Davis AP, et al.
Comparative Toxicogenomics Database (CTD): update 2021.
\textit{Nucleic Acids Res}. 2021;49(D1):D1138-D1143.

\bibitem{xiong2022ddinter}
Xiong G, et al.
DDInter: an online drug-drug interaction database towards improving clinical decision-making and patient safety.
\textit{Nucleic Acids Res}. 2022;50(D1):D1200-D1207.

\end{thebibliography}

\end{document}
